% Source: https://artofproblemsolving.com/community/c3h2900248_2022_jmpsc_d1r1_problem_19
% Stable identity: aops-topic-2900248-q01
The small hexagon's apothem is $\sqrt3$. The tangent side of the large hexagon lies at distance $1$ beyond the midpoint of each small side, so the large apothem is $\sqrt3+1$.
\begin{asy}
usepackage("amsmath");
size(4cm);
defaultpen(fontsize(10pt));
real ap=sqrt(3);pair[] P,Q;for(int i=0;i<6;++i){P[i]=2*dir(60*i);Q[i]=(2+2*sqrt(3)/3)*dir(60*i);}
draw(P[0]--P[1]--P[2]--P[3]--P[4]--P[5]--cycle);draw(Q[0]--Q[1]--Q[2]--Q[3]--Q[4]--Q[5]--cycle);
for(int i=0;i<6;++i){pair M=(P[i]+P[(i+1)%6])/2;real a=degrees(M);draw(arc(M,1,a-90,a+90));}
pair M0=(P[0]+P[1])/2;draw((0,0)--M0,dashed);draw(M0--M0+dir(degrees(M0)),red);dot((0,0));label("$\sqrt3$",M0/2,SE);label("$1$",M0+0.5*dir(degrees(M0)),SE,red);
\end{asy}

A regular hexagon with apothem $h$ has side $\dfrac{2h}{\sqrt3}$. So the side is $\dfrac{2(\sqrt3+1)}{\sqrt3}=\boxed{2+\dfrac{2\sqrt3}3}$.
